\clearpage
\VerdierPhysicalPageAnchor{0170}%
\begingroup\small\setlength{\parskip}{0pt}\setlength{\emergencystretch}{2.5em}
\VerdierSemanticAnchor{ch4-4-2-7-second-diagram}
同じ考察を図式
\[
\begin{tikzcd}[column sep=4.6em,row sep=2.8em]
&H^n(f_2f_1)\arrow[d]\arrow[dr]&\\
H^n(f_1)\arrow[ur]\arrow[r]&
H^n(f_3f_2f_1)\arrow[r]&H^n(f_3f_2),
\end{tikzcd}
\]
に適用すると、同型
\VerdierSemanticAnchor{ch4-4-2-7-2}
\[
E\!\begin{bmatrix}n\\ f_1,f_2,f_3\end{bmatrix}
\xrightarrow{\ \sim\ }
\operatorname{Im}\bigl(H^n(f_2f_1)\longrightarrow H^n(f_3f_2f_1)\bigr)
\Big/
\operatorname{Im}\bigl(H^n(f_1)\longrightarrow H^n(f_3f_2f_1)\bigr).
\tag{4.2.7.2}
\]
を得る。

\VerdierSemanticAnchor{ch4-4-3}
\medskip\noindent\textbf{4.3. 項
$E\!\begin{bmatrix}n\\ -\infty,p,q,+\infty\end{bmatrix}$。}

\VerdierSemanticAnchor{ch4-4-3-1}
\medskip\noindent\textbf{4.3.1.}
$\widetilde{\mathbb Z}$ を、順序集合 $\mathbb Z$ に最小元 $-\infty$ と
最大元 $+\infty$ を付け加えた順序集合とする。記法を濫用して、順序集合
$\widetilde{\mathbb Z}$ に付随する圏も $\widetilde{\mathbb Z}$ で表す。
$\widetilde{\mathbb Z}$ の射とは、したがって組 $p\leq q$ のことである。
ここで $p$ と $q$ は $\widetilde{\mathbb Z}$ の元である。順序付けられた列
$p_1\leq p_2\leq\cdots\leq p_n$ は
$\mathcal{F}\ell_{n-1}(\widetilde{\mathbb Z})$ の対象を定め、これを
$(p_1,p_2,\ldots,p_n)$ と記す。$\widetilde{\mathbb Z}$ の射で、その始域と
終域が $\mathbb Z$ に属するものを\textit{有限}という。
$\mathcal{F}\ell_n(\widetilde{\mathbb Z})$ の対象で、それを構成する射が
有限であるものを\textit{有限}という。

\VerdierSemanticAnchor{ch4-4-3-2}
\medskip\noindent\textbf{4.3.2.}
\textit{本稿では型 $\widetilde{\mathbb Z}$ のスペクトル対象のみを用い、
以後これらを単にスペクトル対象と呼ぶ。}こうして [1] 第 XV 章 §7 の概念を
再び得る。
$(H^n(p,q),\delta,\ p\leq q\in\widetilde{\mathbb Z})$ を、アーベル圏
$\mathcal A$ に値をとるスペクトル対象とする。次のように置く：
\VerdierSemanticAnchor{ch4-4-3-2-1}
\[
F^pH^n(-\infty,+\infty)
=\operatorname{Im}\bigl(H^n(-\infty,p)\longrightarrow
H^n(-\infty,+\infty)\bigr).
\tag{4.3.2.1}
\]
$F^{+\infty}H^n(-\infty,+\infty)=H^n(-\infty,+\infty)$ かつ
$F^{-\infty}H^n(-\infty,+\infty)=0$ である。項
$F^pH^n(-\infty,+\infty)$ は $\widetilde{\mathbb Z}$ 型の
$H^n(-\infty,+\infty)$ 上の増大フィルトレーションをなす。さらに、各組
$p\leq q$ に対して同型 (4.2.7.2)
\VerdierSemanticAnchor{ch4-4-3-2-filtration-quotient}
\[
E\!\begin{bmatrix}n\\ -\infty,p,q,+\infty\end{bmatrix}
\simeq
F^qH^n(-\infty,+\infty)\big/F^pH^n(-\infty,+\infty).
\]
がある。
\endgroup
